\documentclass[journal]{IEEEtran}
\usepackage[spanish]{babel}
\usepackage[utf8]{inputenc}
\usepackage[T1]{fontenc}
\usepackage{amsmath,amssymb,amsfonts}
\usepackage{graphicx}
\usepackage{cite}
\usepackage{url}
\usepackage[hidelinks]{hyperref}
 

\addto\captionsspanish{\renewcommand{\abstractname}{Abstract}}
\renewcommand{\IEEEkeywordsname}{Palabras clave}
 
\begin{document}
 
\title{Esteganografía por Enmascaramiento con Eco}
 
\author{Nasser~Brown, David~Duarte y Jose~Zúñiga%
\thanks{N.Brown, D. Duarte y J. Zuñiga son estudiantes de la Maestría en Electrónica,
Escuela de Ingeniería Electrónica, curso de Procesamiento Digital de Señales
(correo electrónico: nasserbrown@estudiantec.cr; davidduarte@estudiantec.cr; je.zuniga@estudiantec.cr).}}
 
\markboth{Procesamiento Digital de Señales, Maestría en Electrónica}%
{Brown, Duarte y Zúñiga: Esteganografía por Enmascaramiento con Eco}
 
\maketitle
 
\begin{abstract}
This paper presents the design, implementation, and evaluation of the project.
A summary of the problem, the proposed method, the main results, and the
conclusions should be written here in English.
\end{abstract}
 
\begin{IEEEkeywords}
Autocorrelación del cepstro, Convolución, Enmascaramiento con eco,
Esteganografía de audio, Solapamiento y suma.
\end{IEEEkeywords}
 
\section{Introducción}
\IEEEPARstart{E}{n} esta sección se debe describir el problema de investigación,
la motivación del proyecto, los objetivos generales y específicos, así como el
alcance del trabajo. El ocultamiento de información mediante eco fue propuesto
originalmente en \cite{gruhl1996echo,bender1996techniques}.
 
\section{Metodología}
Se describe el enfoque metodológico, la arquitectura del sistema, las
herramientas utilizadas, los datos de entrada y el procedimiento experimental.
 
\section{Resultados y Discusión}
Se presentan los hallazgos principales, las métricas de evaluación y la
comparación con soluciones previas o referentes del estado del arte.
 
\section{Conclusiones}
Se sintetizan los principales logros, limitaciones y propuestas de trabajo
futuro.
 
\bibliographystyle{IEEEtran}
\bibliography{referencias}
 
\end{document}
